\section{Limitations}\label{sec:limitations}

Every limitation below is one that a commercial gas desk has already paid to
remove, which is the immediate disadvantage any public-data study of this market
starts from. They are listed so that the scope of the null is exact.

\paragraph{Terrestrial AIS misses vessels and drops them.} The archives and the
live feed see a vessel only within range of a shore receiver. Vessels that load out
of range leave no departure, vessels at sea are invisible for the whole crossing,
and the live feed additionally loses any vessel that does not hold a subscription
slot. The rescue backstop of Section~\ref{sec:live} quantifies the live shortfall:
\nRescueSkipped{} candidate lookups were declined by the budget in the period
shown, and \nRescueLookups{} were made. The measured publication lag of a
departure is negligible on the archives and up to a day on the live feed.

\paragraph{The polygons are hand-drawn.} The \nZonePolygons{} berth, anchorage and
approach polygons were digitised by hand and are provisional at several terminals.
An envelope drawn too tight emits spurious exits during a channel transit; one
drawn too loose merges neighbouring terminals. The state machine's stickiness and
berth-override rules limit the damage, and the capture check bounds the aggregate
effect on the US side, but no per-terminal audit of the geometry has been done.

\paragraph{The two archives have opposite capture gradients.} NOAA-only capture of
US export cargoes is \captureTwentyTwenty{} in 2020 and \captureTwentyEighteen{} in
2018 (Table~\ref{tab:capture}), and GFW's capture runs the other way. Early-year
levels of every US signal are biased low. The walk-forward designs operate on
within-window changes and the FWL scan on residuals, but a level-dependent
non-linearity in the early years would be misread.

\paragraph{No European anchorage history.} The GFW feed carries port visits but no
anchorage events, so European queue and absorption signals exist only from the 2026
live feed. Four of the five composite signals originally pre-registered as features
for Section~\ref{sec:partb} turned out to have between 9 and 46 daily observations
and were excluded by a documented amendment before any spread fit. The at-sea and
US-side mechanism was tested on the decade; the EU-congestion mechanism could not be
tested.

\paragraph{Destination resolution.} Between 75\% and 83\% of the at-sea stock has no
resolved destination. The EU-bound series is a minority observed slice whose
definition depends on coverage. The unresolved band is carried as its own feature,
and the ratio of the two was tested as H2. Attenuation from this measurement error
biases coefficients toward zero and cannot be excluded as a contributing factor.

\paragraph{The coverage seam.} The backfill feeds end on \noaaEnds{} and \gfwEnds.
The full-panel results in Appendix~\ref{app:fullpanel} include the thin live tail.
In that sample 24\% of the holdout weeks lie after the seam, and one full-panel
finding, a change of sign in the H2 holdout coefficient, did not survive truncation
and is not claimed. The primary sample is truncated at the seam.

\paragraph{Regime pooling.} Signals are read from the pooled regime tag, which is
the only tag spanning the decade. A small residual over-count where NOAA and GFW both
observe the same US departure is documented in the project's signal catalogue. A
robustness run on the NOAA tag alone for the US-side signals was not performed.

\paragraph{The pre-registration trail.} As Section~\ref{sec:prereg} states, the git
history does not independently establish that the log entries preceded the fits they
govern.

\paragraph{Sample size and a dominant regime.} The scored spread sample is
\bWeeksCov{} weeks, heavily autocorrelated, and 2022 alone spans a range larger than
the remainder of the decade combined. The year-consistency condition and the holdout
are the defences against this; a longer sample without a coverage seam would be
preferable to either.

\paragraph{The outage monitor is slow.} The usable detector has a twelve-day median
delay. A terminal outage of that length is public news well before the monitor
fires, so its value is as an automated check on the reconstruction, not as
information.

\paragraph{Vintage residual.} \vintagePostDay{} rows in the live tail changed after a
next-day print (Section~\ref{sec:bounds}). They lie outside the modelling sample, but
they show that the knowable bound is exact on event time and only approximately
exact on observation time in a thin-coverage regime.
